\documentclass[11pt]{article}
\usepackage[utf8]{inputenc}
\usepackage[T1]{fontenc}
\usepackage{amsmath,amssymb,mathtools}
\usepackage{geometry}
\usepackage{booktabs}
\geometry{margin=1in}

\title{Paired Timestep-Shortcut Training Objective}
\author{}
\date{}

\begin{document}
\maketitle

\section{Active configuration}

This document describes the intended training path for branch
\texttt{feat/timestep-shortcut-output-distill} after the paired-batch update.
The default timestep-shortcut setup is:
\begin{center}
\begin{tabular}{ll}
\toprule
Switch & Value \\
\midrule
\texttt{--use-timestep-shortcut} & on \\
\texttt{--timestep-paired-batch} & on \\
\texttt{--timestep-shortcut-layer} & $k=9$ by default \\
\texttt{--timestep-shortcut-loss-mode} & \texttt{activation\_huber} \\
\texttt{--timestep-shortcut-use-layer-cond} & off by default \\
\texttt{--shortcut-predictor} & \texttt{hybrid\_depth10} alias for \texttt{hybrid\_deep\_10} \\
\texttt{--no-shortcut-predictor-normalize-input} & predictor receives raw layer-$k$ activations \\
\texttt{--use-timestep-direct-loss} & on \\
\texttt{--use-timestep-bootstrap-loss} & on \\
\texttt{--use-timestep-output-distill} & on \\
\texttt{--no-use-depth-time-resync} & depth-time losses inactive \\
\texttt{--no-output-distill} & legacy depth output-distill inactive \\
\texttt{--private-loss} & off by default; cosine-squared private loss inactive \\
\bottomrule
\end{tabular}
\end{center}

The direct and bootstrap timestep losses are Smooth-L1/Huber activation losses.
They are not direction--magnitude losses in this mode.

\section{Paired local batch}

Let the global batch size be $B=128$. With paired timestep training, only
$B/2=64$ distinct images are used per global step. Each image is materialized at
two noise levels inside the same local device batch. If there are $D$ devices,
the local batch size is $B_{\mathrm{loc}}=B/D$ and the local number of distinct
images is
\[
  M_{\mathrm{loc}} = \frac{B_{\mathrm{loc}}}{2}.
\]
The implementation requires $B_{\mathrm{loc}}$ to be even.

For local image $i\in\{1,\ldots,M_{\mathrm{loc}}\}$, draw one clean latent
$x_{0,i}$, one Gaussian noise sample $x_{1,i}\sim\mathcal{N}(0,I)$, and three
ordered timestep indices
\[
  q_a < q_b < q_c.
\]
The repo convention is that $\tau=0$ is pure noise and larger $\tau$ is cleaner.
Thus $a$ is the heavier-noise state, $b$ is lighter, and $c$ is the lightest.
With $T=50$ shortcut bins,
\[
  \tau_a=\frac{q_a}{T-1},\qquad
  \tau_b=\frac{q_b}{T-1},\qquad
  \tau_c=\frac{q_c}{T-1}.
\]
The noisy latents are
\[
  x_a=(1-\tau_a)x_1+\tau_a x_0,\qquad
  x_b=(1-\tau_b)x_1+\tau_b x_0,\qquad
  x_c=(1-\tau_c)x_1+\tau_c x_0.
\]
Only $x_a$ and $x_b$ are put through the full DiT backbone. $x_c$ is used only
to form cheap input/timestep embeddings for the bootstrap target condition.

\section{Backbone forward passes}

The heavy and light states are concatenated inside each local batch:
\[
  X_{\mathrm{pair}} = [x_a; x_b],\qquad
  \Tau_{\mathrm{pair}} = [\tau_a; \tau_b],\qquad
  Y_{\mathrm{pair}}=[y;y].
\]
A single DiT apply over this paired batch is equivalent in cost to two half-batch
backbone forwards:
\[
  \bigl(\hat v_a,\{Z_\ell^a\}_{\ell=0}^{L},e_a\bigr),
  \bigl(\hat v_b,\{Z_\ell^b\}_{\ell=0}^{L},e_b\bigr)
  =
  F_\theta^{\mathrm{hidden}}([x_a;x_b],[\tau_a;\tau_b],[y;y]).
\]
Here $e_a,e_b$ are DiT timestep embeddings, $Z_0$ is the patch embedding, and
$Z_k$ is the output of block/layer $k$. The default fixed shortcut layer is
$k=9$.

The third timestep is not fully forwarded through DiT. The code computes only
\[
  (Z_0^c,e_c)=F_\theta^{\mathrm{input\_emb}}(x_c,\tau_c,y),
\]
which gives the target patch embedding and timestep embedding needed by the
predictor.

\section{Base velocity loss}

The normal flow-matching velocity target is independent of timestep:
\[
  v^\star=x_0-x_1.
\]
The base loss is evaluated on both paired noise levels:
\[
  \mathcal{L}_{\mathrm{gen}}
  =
  \frac{1}{2M_{\mathrm{loc}}}
  \sum_i
  \left(
    \left\lVert \hat v_{a,i}-v_i^\star\right\rVert_2^2
    +
    \left\lVert \hat v_{b,i}-v_i^\star\right\rVert_2^2
  \right).
\]

\section{Predictor interface}

The timestep shortcut predictor is
\[
  P_\phi(S,e_{\mathrm{src}},e_{\mathrm{tgt}},Z_0^{\mathrm{tgt}},q_{\mathrm{src}},q_{\mathrm{tgt}})
  \mapsto \widehat Z.
\]
For this paired timestep mode:
\begin{itemize}
\item $S$ is the raw source activation, not L2-normalized.
\item the layer is fixed to $k$ globally, so source/target layer ids are not
  active predictor conditions by default;
\item the predictor still receives source and target timestep information
  $(e_{\mathrm{src}},e_{\mathrm{tgt}},q_{\mathrm{src}},q_{\mathrm{tgt}})$;
\item the target patch embedding $Z_0^{\mathrm{tgt}}$ can condition the
  predictor.
\end{itemize}

\section{Smooth-L1 activation loss}

For a predicted activation $\widehat Z$ and a target activation $Z$, define the
per-sample RMS normalizer
\[
  r(Z)_i =
  \sqrt{\mathrm{mean}_{n,d}\left(Z_{i,n,d}^2\right)+\epsilon}.
\]
The normalized residual is
\[
  \rho_i=\frac{\widehat Z_i-\mathrm{sg}(Z_i)}{\mathrm{sg}(r(Z)_i)}.
\]
The Smooth-L1/Huber activation loss is
\[
  \mathcal{H}_\delta(\widehat Z,Z)
  =
  \mathrm{mean}_{i,n,d}\left[
  \begin{cases}
    \frac{1}{2}\rho_{i,n,d}^2, & |\rho_{i,n,d}|\le \delta,\\
    \delta\left(|\rho_{i,n,d}|-\frac{1}{2}\delta\right),
      & |\rho_{i,n,d}|>\delta,
  \end{cases}
  \right],
\]
with default $\delta=1$.

\section{Direct timestep loss}

The heavy activation at layer $k$ predicts the light activation at the same
layer:
\[
  \widehat Z_k^b
  =
  P_\phi(Z_k^a,e_a,e_b,Z_0^b,q_a,q_b).
\]
The direct loss is
\[
  \mathcal{L}_{\mathrm{direct}}
  =
  \lambda_{\mathrm{direct}}
  \mathcal{H}_\delta(\widehat Z_k^b,Z_k^b),
\]
where the implementation uses \texttt{--shortcut-lambda-dir} as
$\lambda_{\mathrm{direct}}$.

\section{Bootstrap loss}

The online two-hop prediction is
\[
  \widehat Z_k^{c,\mathrm{2hop}}
  =
  P_\phi(\widehat Z_k^b,e_b,e_c,Z_0^c,q_b,q_c).
\]
The EMA one-hop target is
\[
  \widehat Z_k^{c,\mathrm{ema}}
  =
  P_{\bar\phi}(\mathrm{sg}(Z_k^a),\mathrm{sg}(e_a),e_c,Z_0^c,q_a,q_c).
\]
The bootstrap equation is therefore
\[
  P_\phi(P_\phi(a,b),c)\approx P_{\bar\phi}(a,c).
\]
The bootstrap loss is
\[
  \mathcal{L}_{\mathrm{boot}}
  =
  \lambda_{\mathrm{boot}}
  \mathcal{H}_\delta
  \left(
    \widehat Z_k^{c,\mathrm{2hop}},
    \mathrm{sg}(\widehat Z_k^{c,\mathrm{ema}})
  \right),
\]
where $\bar\phi$ is the predictor EMA.

\section{Output distillation}

The same direct prediction $\widehat Z_k^b$ is reused for output distillation.
It is fed into the DiT tail from layer $k+1$ onward, but now with the light
timestep and class condition:
\[
  \widehat v_b^{\mathrm{tail}}
  =
  F_\theta^{\mathrm{tail}}
  (x_b,\tau_b,y;\ \mathrm{resume\_hidden}=\widehat Z_k^b,\ 
   \mathrm{resume\_start}=k+1).
\]
The target is the full light-noise DiT prediction from the paired backbone
forward:
\[
  v_b^{\mathrm{target}}=\mathrm{sg}(\hat v_b).
\]
The output distillation loss is the relative MSE
\[
  \mathcal{L}_{\mathrm{out}}
  =
  \frac{1}{M_{\mathrm{loc}}}
  \sum_i
  \frac{
    \left\lVert \widehat v_{b,i}^{\mathrm{tail}}
      - v_{b,i}^{\mathrm{target}}\right\rVert_2^2
  }{
    \left\lVert v_{b,i}^{\mathrm{target}}\right\rVert_2^2 + 10^{-6}
  }.
\]
Unless \texttt{--timestep-output-distill-train-tail} is enabled, this branch
updates the predictor while the DiT tail parameters are stop-gradient.

\section{Inactive cosine-squared private loss}

The cosine-squared private activation loss is off by default in this run. In
command terms, do not pass \texttt{--private-loss}. Thus
\[
  \mathcal{L}_{\mathrm{private}}=0.
\]

\section{Total loss}

With depth-time resync, legacy depth shortcut losses, legacy output distill,
skip-FM, and cosine-squared private loss inactive, the optimized loss is
\[
\boxed{
  \mathcal{L}_{\mathrm{total}}
  =
  \mathcal{L}_{\mathrm{gen}}
  +
  \mathcal{L}_{\mathrm{direct}}
  +
  \mathcal{L}_{\mathrm{boot}}
  +
  \lambda_{\mathrm{out}}\mathcal{L}_{\mathrm{out}}.
}
\]
For the recommended command, $\lambda_{\mathrm{out}}=0.05$.

\section{Compute summary}

Per training step, the timestep losses require:
\begin{itemize}
\item one paired DiT forward over $[x_a;x_b]$, equivalent to two half-batch
  backbone forwards and no larger than the original global batch cost;
\item one cheap input/timestep embedding pass for $x_c$;
\item three predictor calls:
  \[
    P_\phi(a,b),\qquad P_\phi(P_\phi(a,b),c),\qquad P_{\bar\phi}(a,c).
  \]
\end{itemize}

\section{FID evaluation variants}

Evaluation now reports three 4k-sample FID variants when
\texttt{--num-fid-samples 4096}.

\paragraph{1. Full DiT FID.}
The standard sampler uses the full DiT at every denoise step:
\[
  \texttt{val/FID},\qquad \texttt{val/sFID}.
\]

\paragraph{2. One-step chained timestep shortcut FID.}
At each denoise step $j$, the latent is run through the DiT prefix up to layer
$k=9$, then the predictor maps the layer-$9$ activation from step $j$ to
step $j+1$. The DiT tail then runs with the timestep embedding and class
condition of step $j+1$:
\[
  Z_k^{j+1}=P_\phi(Z_k^j,e_j,e_{j+1},Z_0^{j+1},q_j,q_{j+1}),
  \qquad
  \hat v_{j+1}=F_\theta^{\mathrm{tail}}(Z_k^{j+1},e_{j+1},y).
\]
This is evaluated as
\[
  \texttt{val/FID\_timestep\_chain},\qquad
  \texttt{val/sFID\_timestep\_chain}.
\]
The implementation avoids a full future-backbone pass; it only computes the
target timestep/input embeddings before the predictor.

\paragraph{3. One-prefix fanout timestep shortcut FID.}
The fanout sampler runs the DiT prefix once at the first denoise step, then
rolls the layer-$9$ activation forward through the predictor across the
remaining reference timesteps:
\[
  Z_k^{j+1}=P_\phi(Z_k^j,e_j,e_{j+1},Z_0^{j+1},q_j,q_{j+1}),
  \qquad j=0,\ldots,48.
\]
Each predicted $Z_k^{j+1}$ is passed through the DiT tail with the corresponding
timestep embedding. To avoid OOM on TPU v5p-8, the predictor/tail rollout is
implemented with \texttt{jax.lax.scan}; it does not materialize all 49 tail
activations at once. This is evaluated as
\[
  \texttt{val/FID\_timestep\_fanout},\qquad
  \texttt{val/sFID\_timestep\_fanout}.
\]

The full sampler still uses
\[
  \texttt{--fid-num-steps}=50.
\]
The older two-step skip-25 ablation remains available but is off by default:
\[
  \texttt{--fid-timestep-skip-num-steps}=2,\qquad
  \texttt{--fid-timestep-skip-jump}=25,\qquad
  \texttt{--fid-timestep-skip-every}=1.
\]

\end{document}
